% ##################################################################################################################
\chapter{Padang}
\label{ch:padang}
\hfill \textbf{Author:} Gregor Lämmel

\editdone{This text has undergone the professional edit. Please no grammatical changes anymore! They are most-probably wrong.}

% ##################################################################################################################
The Padang scenario demonstrates the \gls{matsim} application to large-scale evacuation problems. The scenario has been created as part of the third party funded project, ``Last-Mile''. \citet{00TaubenboeckEtAl2012ConcludingLastMilePaperNatHazards} gives a comprehensive overview.
Padang is located on the west coast of Sumatra Island, Indonesia. In 2014, the city had a population of about 1\,000\,000\,people. 
Because of its problematic location on the coast in a so-called ``seismically locked'' area~\citep{McCloskey2010Padang2009Earthquake}, Padang is prone to earthquakes and subsequent tsunamis. In the ``Last-Mile'' project, a realistic tsunami scenario, triggered by an earthquake about 300\,km off the coast, was identified~\citep{GosebergSchlurmann2009HazardMappingPadang}. The assumed tsunami would leave about 30\,minutes for the evacuation. The flooding would reach as far as three kilometers inland, thus threatening up to 330\,000\,lives. \citet{Laemmel_PhDThesis_2011} developed a \gls{matsim} scenario representing the city with its affected population. One unusual aspect of the Padang situation is the expected universal evacuation by foot; simulating pedestrians with \gls{matsim} was a novelty when this project started. The standard simulation model (see, \eg Chapter~\ref{sec:trafficflowmodel}) was thus adapted to deal with pedestrians. 
Details are discussed in \citet{00LaemmelKluepfelNagel2009EvacPadangAtBookTimmermanns}. Another important variation, contrary to most standard transport scenarios, is that network links would flood once the tsunami reached them. Thus, accessibility---flooded or not flooded---of the network links is time-dependent, which is modeled by a time-dependent network \citep{LaemmelGretherNagel2009TimeDependentNetworks}. In the time-dependent network concept, link attributes---like \emph{freespeed}---can be changed, while the simulation is running, by precomputed network change events. For the Padang scenario, the network change events have been extracted from microscopic flooding simulation data.

% ##################################################################################################################
%
\createfigure%
{Evacuation time analysis for downtown Padang.}%
{Evacuation time analysis for downtown Padang. Numbers showing average evacuation time in minutes, which are also indicated by the colors green, yellow, red}%
{\label{chap:using:padang}}%
{\includegraphics[width=0.6\textwidth, angle=0]{scenarios/figures/dwntwnpdg}}%
{}

Key Padang scenario facts:
\begin{itemize}\styleItemize
\item The network consists of about 6\,000\,nodes and 17\,000\,links.
\item Synthetic populations for morning, afternoon, and night have been created, containing up to 330\,000\,agents.
\item The flooding is modeled by a set of 109\,network change events (one per minute), affecting 7\,609\,links.
\item A set of 42\,shelter buildings, which could be used for vertical evacuation, is also part of the scenario.
\end{itemize}
Based on the Padang scenario, various evacuation strategies have been investigated:
\begin{itemize}\styleItemize
\item A seemingly obvious evacuation strategy is the shortest path solution, where everyone is on the shortest path. This solution, however, ignores possible congestion and lead to unfeasible results;
\item Shorter evacuation times are achieved with a Nash equilibrium approach, where everyone tries to find an optimal evacuation route through iterative learning \citep{00LaemmelKluepfelNagel2009EvacPadangAtBookTimmermanns};
\item While the Nash equilibrium reduces individual evacuation time, total evacuation time might not be minimal. The marginal social cost-based simulation approach tries to minimize the total evacuation time \citep{00LaemmelFloetteroed2009KISysOptEvac,00DresslerFloetteroedLaemmelNagelSkutella2010OptimalEvacuationLargeScaleScenarios};
\item These three basic evacuation approaches are investigated in combination of flooding \citep{LaemmelGretherNagel2009TimeDependentNetworks,Laemmel_PhDThesis_2011};
\item Further, an evacuation strategy to reduce the exposure to risk has been developed by \citep{00LaemmelKluepfelNagel2010PEDRiskPrinted};
\item And finally, \citet{00FloetteroedLaemmel2010ICECShelterEvac} proposes a method to integrate shelter buildings, which are evacuation sinks (\ie safe places) 
% "\Karen{ What?? sink? Does he mean 'shed', or some kind of building? Or do the buildings have sanitary facilities? Need help with this word.}" 
with limited capacity, into the simulation.  
\end{itemize}

% ##################################################################################################################
